%% Работа с русским языком
\usepackage{cmap}			 % поиск в PDF
\usepackage{mathtext} 		 % русские буквы в формулах
\usepackage[T2A]{fontenc}	 % кодировка
\usepackage[utf8]{inputenc}	 % кодировка исходного текста
\usepackage[russian]{babel}	 % локализация и переносы

%% Пакеты для работы с математикой
\usepackage{amsmath,amsfonts,amssymb,amsthm,mathtools}
\usepackage{icomma}

%% Нумерация формул (опционально)
%\mathtoolsset{showonlyrefs=true} % показывать номера только у тех формул, на которые есть \eqref{} в тексте.
%\usepackage{leqno}               % нумерация формул слева

%% Шрифты
\usepackage{euscript}	 % шрифт "Евклид"
\usepackage{mathrsfs}    % красивый мат. шрифт

%% Некоторые полезные макросы для дебага
\makeatletter
\newcommand\thefontsize{The current font size is: \f@size pt} % пример: \section{\thefontsize}
\makeatother

%% Настройка заголовков
\usepackage{titlesec}

\titleformat{\section}
    {\normalfont\bfseries\centering\fontsize{14}{17}\selectfont}
    {\thesection}{1em}{}

\titleformat{\subsection}
    {\normalfont\bfseries\centering\fontsize{14}{17}\selectfont}
    {\thesubsection}{1em}{}

\titleformat{\subsubsection}
    {\normalfont\bfseries\centering\fontsize{14}{17}\selectfont}
    {\thesubsubsection}{1em}{}

\titlespacing*{\section}{0pt}{1.5em}{1em}
\titlespacing*{\subsection}{0pt}{1.2em}{0.8em}
\titlespacing*{\subsubsection}{0pt}{1em}{0.6em}

%% Поля (геометрия страницы)
\usepackage[left=3cm,right=1.5cm,top=2cm,bottom=2cm,bindingoffset=0cm]{geometry}

%% Русские списки
\usepackage{enumitem}
\makeatletter
\AddEnumerateCounter{\asbuk}{\russian@alph}{щ}
\makeatother

%% Работа с картинками
\usepackage{caption}
\usepackage{subcaption}
\usepackage{graphicx}
\graphicspath{{images/}{images2/}}
\setlength\fboxsep{3pt}
\setlength\fboxrule{1pt}
\usepackage{wrapfig}
\addto\captionsrussian{
    \renewcommand{\figurename}{Рисунок}
    \renewcommand{\tablename}{Таблица}
}
\DeclareCaptionLabelSeparator{gostdash}{~--~}
\captionsetup{
    labelsep=gostdash
}
\captionsetup[figure]{
    justification=centering,
    singlelinecheck=true
}

%% Работа с таблицами
\usepackage{array,tabularx,tabulary,booktabs}
\usepackage{longtable}
\usepackage{multirow}
\usepackage{makecell}
\usepackage{hhline}
\usepackage{siunitx}
\captionsetup[table]{
    justification=raggedright,
    singlelinecheck=false,
    position=top
}

%% Красная строка
\setlength{\parindent}{1.25cm}

%% Интервалы
\usepackage{setspace}
\onehalfspacing

%% TikZ
\usepackage{tikz}
\usetikzlibrary{graphs,graphs.standard}

%% Колонтитулы
\usepackage{fancyhdr}

\pagestyle{fancy}
\fancyhf{}
\fancyfoot[C]{\thepage}

\renewcommand{\headrulewidth}{0pt}
\renewcommand{\footrulewidth}{0pt}

\fancypagestyle{plain}{
    \fancyhf{}
    \fancyfoot[C]{\thepage}
    \renewcommand{\headrulewidth}{0pt}
    \renewcommand{\footrulewidth}{0pt}
}

%% Перенос знаков в формулах (по Львовскому)
\newcommand*{\hm}[1]{#1\nobreak\discretionary{}{\hbox{$\mathsurround=0pt #1$}}{}}

%% Дополнительно
\usepackage{float}
\usepackage{gensymb}
\usepackage{listings}
\lstset{
basicstyle=\small\ttfamily,
columns=flexible,
breaklines=true
}

%% Hyperref (для ссылок внутри pdf)
\usepackage[unicode, pdftex]{hyperref}

%% Отступ перед первым абзацем в каждом разделе
\usepackage{indentfirst}
